\chapter{Propuesta}\label{chapter:proposal}

En este capítulo se presenta BT-GRAPHRAG, la propuesta de solución al problema abordado en la presente investigación, en correspondencia con los objetivos previamente definidos para la misma. Se describe su concepción general, los fundamentos teóricos que la sustentan y las decisiones de diseño adoptadas, así como los principales componentes que la integran y su articulación dentro del \emph{pipeline} de GraphRAG. Asimismo, se exponen los mecanismos incorporados para abordar las limitaciones identificadas en el estado del arte, con énfasis en el tratamiento explícito de la temporalidad en grafos de conocimiento.

La propuesta responde a las seis limitaciones identificadas en la sección~\ref{sec:soa-conclusions} y sintetizadas en la Tabla~\ref{tab:soa-comparison}. Su principio fundamental consiste en incorporar la temporalidad como una dimensión de primera clase dentro del \emph{pipeline} de GraphRAG. Esta integración no implica la sustitución de los componentes ya existentes en GraphRAG~\cite{edge2024local}; por el contrario, se preservan la sumarización jerárquica basada en comunidades de Leiden~\cite{traag2019leiden} y los mecanismos de búsqueda por comunidades. Sobre esta base, se introducen cuatro mecanismos adicionales:
\begin{itemize}
    \item un modelo bitemporal a nivel de arista,
    \item una clasificación de cardinalidad asimétrica,
    \item un mecanismo de resolución temporal de entidades,
    \item un módulo de detección y resolución activa de conflictos durante la indexación.
\end{itemize}

El capítulo se organiza en ocho secciones. La sección~\ref{sec:bt-overview} presenta los objetivos de diseño y establece la correspondencia entre cada componente propuesto y las limitaciones identificadas en el estado del arte. La sección~\ref{sec:bt-pipeline} describe la arquitectura del \emph{pipeline} de indexación y sus siete etapas. En la sección~\ref{sec:bt-bitemporal} se formaliza la cuádrupla de estado temporal y los estados epistémicos derivados de esta. La sección~\ref{sec:bt-cardinality} introduce la clasificación de cardinalidad en cuatro categorías y su papel en la orientación de la detección de conflictos. La sección~\ref{sec:bt-cger} aborda los mecanismos de resolución de entidades y de tipos de relación, delegando en el LLM la decisión final de fusión a partir de las descripciones y los períodos activos de las entidades candidatas. La sección~\ref{sec:bt-etcdr} describe el módulo de detección bidireccional de conflictos y el componente encargado de su resolución. La sección~\ref{sec:bt-integration} expone la integración con la pirámide de comunidades, la descomposición temporal de consultas y la actualización incremental del corpus. Finalmente, la sección~\ref{sec:bt-conclusions} presenta la correspondencia uno a uno entre las limitaciones del estado del arte y los componentes propuestos para su solución.

%% ============================================================
\section{Visión general de BT-GRAPHRAG}\label{sec:bt-overview}

BT-GRAPHRAG (\emph{Bitemporal GraphRAG}) es un sistema de Generación Aumentada por Recuperación basado en grafos de conocimiento. Su rasgo distintivo es que cada arista del grafo incorpora, además de la tripleta sujeto-relación-objeto, una marca temporal estructurada que registra tanto cuándo el hecho fue verdadero en el mundo real como cuándo el sistema llegó a creerlo.

La propuesta se articula en torno a cinco objetivos fundamentales de diseño:

\begin{itemize}
    \item \textbf{Bitemporalidad nativa.} Cada arista almacena una cuádrupla de marcas temporales que separa el tiempo válido del tiempo de transacción. Esta separación se fundamenta en el marco formal de TLINQ y VLINQ~\cite{tlinq2022}, y su novedad reside en que se integra en el propio \emph{pipeline} de GraphRAG, en lugar de delegarse a un sistema de memoria externo como ocurre en ZEP/Graphiti~\cite{zep2025}.
    \item \textbf{Resolución activa de conflictos.} Los sistemas T-GRAG~\cite{tgrag2025} y TG-RAG~\cite{tg_rag2025} delegan la coherencia temporal al filtrado en tiempo de consulta. BT-GRAPHRAG, en cambio, detecta los conflictos durante la fase de indexación y aplica en ese mismo momento la estrategia de resolución correspondiente, actualizando directamente las marcas temporales de la arista afectada.
    \item \textbf{Asimetría estructural de la cardinalidad.} La etiqueta binaria exclusiva/no-exclusiva empleada en TREK~\cite{trek2025} se sustituye por una clasificación de cuatro categorías que identifica desde qué extremo de la relación opera la restricción de exclusividad. Esta distinción es la que hace posible una búsqueda dirigida y precisa de conflictos temporales.
    \item \textbf{Resolución temporal de entidades.} Para determinar si dos menciones se refieren a la misma entidad, el sistema combina un pre-filtro por similitud coseno entre los \emph{embeddings} de descripción con una verificación por LLM que recibe los períodos activos como parte del contexto. El LLM emite uno de tres veredictos —\texttt{SAME}, \texttt{DIFFERENT\_ENTITY} o \texttt{DIFFERENT\_TEMPORAL}—, lo que permite distinguir entre referentes distintos y estados temporalmente separados de un mismo referente. Este mecanismo evita tanto la fusión inadecuada de estados muy distintos de una misma entidad como la fragmentación errónea de una entidad cuyo nombre haya cambiado a lo largo del tiempo.
    \item \textbf{Integración no invasiva con GraphRAG.} El \emph{pipeline} estándar de GraphRAG~\cite{edge2024local} ---extracción, deduplicación, particionamiento de Leiden, sumarización jerárquica, y búsqueda local y global--- se conserva íntegramente. Los nuevos componentes se incorporan como etapas adicionales y respetan los formatos de salida del sistema base, sin alterar su lógica interna.
\end{itemize}

% TODO: figura fig:bt-graphrag-overview que ilustre los componentes principales de BT-GRAPHRAG y su correspondencia con las limitaciones del estado del arte.

La Tabla~\ref{tab:limitations-to-components} resume cómo cada componente de la propuesta responde a las limitaciones del estado del arte. La correspondencia es directa: a cada limitación se le asigna un componente y la sección del capítulo donde se desarrolla.

\begin{table}[h]
\centering
\small
\begin{tabular}{p{0.45\linewidth}p{0.30\linewidth}p{0.18\linewidth}}
\hline
\textbf{Limitación del SoA} & \textbf{Componente BT-GRAPHRAG} & \textbf{Sección} \\ \hline
1: Temporalidad no integrada al \emph{pipeline} & \emph{Pipeline} de indexación temporal y consulta temporal & \ref{sec:bt-pipeline}, \ref{sec:bt-integration} \\
2: Sistemas evaden la resolución de conflictos & Detección bidireccional y selección de estrategia & \ref{sec:bt-etcdr} \\
3: Cardinalidad simétrica (TREK) & Clasificación en cuatro categorías & \ref{sec:bt-cardinality} \\
4: Bitemporalidad fuera del \emph{pipeline} GraphRAG & Cuádrupla de estado temporal en cada arista & \ref{sec:bt-bitemporal} \\
5: Resolución de entidades sin temporalidad & Verificación por LLM con contexto temporal y veredicto \texttt{DIFFERENT\_TEMPORAL} & \ref{sec:bt-cger} \\
6: Corrección y evolución sin distinción formal & Estados epistémicos derivados de la cuádrupla & \ref{sec:bt-bitemporal} \\ \hline
\end{tabular}
\caption{Correspondencia entre las limitaciones identificadas en el estado del arte y los componentes de BT-GRAPHRAG que las resuelven.}
\label{tab:limitations-to-components}
\end{table}

El alcance de este capítulo se circunscribe al diseño conceptual: el modelo de datos, los algoritmos y la arquitectura del sistema. Los aspectos relativos a tecnologías concretas, configuración e implementación, así como el plan de evaluación, se tratan en capítulos posteriores del presente documento.

%% ============================================================
\section{Arquitectura del pipeline de indexación}\label{sec:bt-pipeline}

El \emph{pipeline} de indexación de BT-GRAPHRAG se organiza en siete etapas secuenciales. Las cuatro primeras procesan cada lote de documentos nuevos antes de modificar el grafo; las tres restantes persisten el grafo y mantienen las estructuras que se construyen a partir de él (la jerarquía de comunidades de Leiden y sus reportes), antes de dar paso a la fase de consulta.

\begin{enumerate}
    \item \textbf{Extracción temporal.} Un LLM procesa cada bloque de texto y extrae entidades y relaciones, de forma análoga a GraphRAG. La diferencia es que cada relación incorpora dos campos temporales adicionales: el instante de inicio y el de fin del hecho descrito. El \emph{prompt} recibe la fecha del documento como contexto y normaliza las expresiones temporales relativas (es decir, convierte expresiones como ``el año pasado'' o ``hace tres meses'' en fechas absolutas tomando la fecha del documento como referencia)~\cite{TODO_temporal_normalization}. Cuando dicha fecha es significativamente posterior a la del hecho, el documento se marca como \emph{late arrival} y activa la lógica \emph{as-of} (que reconstruye el estado del grafo tal como era en la fecha del hecho, evitando que correcciones posteriores contaminen la indexación), descrita en detalle en la sección~\ref{sec:bt-bitemporal}.
    \item \textbf{Resolución contextual de entidades.} Las entidades extraídas se contrastan con las ya presentes en el grafo mediante un proceso en dos fases: un pre-filtro por similitud coseno entre los \emph{embeddings} de descripción y, para los pares cuyo coseno supera el umbral, una verificación por LLM que incorpora el período activo de cada entidad como parte del contexto y devuelve uno de los veredictos \texttt{SAME}, \texttt{DIFFERENT\_ENTITY} o \texttt{DIFFERENT\_TEMPORAL}. El procedimiento completo se detalla en la sección~\ref{sec:bt-cger}.
    \item \textbf{Resolución de tipos de relación.} El mismo principio se aplica a los tipos de relación. El objetivo es consolidar bajo un tipo canónico único aquellas etiquetas semánticamente equivalentes ---por ejemplo, \texttt{IS\_CEO\_OF} y \texttt{LEADS}--- antes de proceder a la evaluación de cardinalidad y a la detección de conflictos.
    \item \textbf{Clasificación de cardinalidad.} Cada tipo de relación se clasifica una única vez en una de las cuatro categorías de cardinalidad definidas en la propuesta. El resultado se almacena como atributo del tipo y determina la dirección de la búsqueda de conflictos. Los detalles se exponen en la sección~\ref{sec:bt-cardinality}.
    \item \textbf{Detección y resolución de conflictos.} Para cada arista candidata, el módulo lanza consultas dirigidas por la cardinalidad del tipo de relación correspondiente (la categoría que indica desde qué extremo de la relación opera la restricción de exclusividad; véase la sección~\ref{sec:bt-cardinality}). Ante la detección de un conflicto, un componente basado en LLM selecciona una de cinco estrategias de resolución y la aplica sobre las marcas temporales de la arista afectada. Este módulo se describe en profundidad en la sección~\ref{sec:bt-etcdr}.
    \item \textbf{Escritura en el almacén bitemporal.} Las aristas y entidades resultantes se almacenan en dos representaciones complementarias y sincronizadas. La primera es una representación tabular compatible con los \emph{artifacts} que produce GraphRAG, lo que permite que la sumarización por comunidades y las búsquedas local y global del sistema base operen sin modificación. La segunda es una representación en grafo que conserva la cuádrupla temporal completa y soporta las consultas dirigidas por cardinalidad. La elección de tecnologías concretas para cada una se aborda en capítulos posteriores.
    \item \textbf{Reportes de comunidad temporales y consulta.} Sobre el grafo se ejecuta el algoritmo de Leiden de GraphRAG~\cite{traag2019leiden}, que genera la jerarquía de comunidades. Cada reporte se enriquece con una narrativa temporal que sintetiza la evolución del subgrafo durante el período correspondiente. En la fase de consulta, las preguntas con múltiples referencias temporales se descomponen en sub-consultas independientes, cuyos resultados se agregan mediante una síntesis temporal. El mecanismo completo se describe en la sección~\ref{sec:bt-integration}.
\end{enumerate}

% TODO: figura fig:bt-pipeline que muestre las siete etapas en secuencia, con sus entradas y salidas.

La decisión arquitectónica central es mantener estas dos representaciones complementarias del grafo. La representación tabular preserva la compatibilidad con las herramientas de GraphRAG; la representación en grafo conserva la cuádrupla temporal completa, requisito indispensable tanto para las consultas \emph{as-of} como para la búsqueda bidireccional del módulo de conflictos. Prescindir de la primera comprometería la naturaleza de BT-GRAPHRAG como extensión de GraphRAG; prescindir de la segunda haría inviable la formulación de dichas consultas.

El esquema de datos amplía las tablas de entidades y relaciones del modelo de GraphRAG con las columnas recogidas en la Tabla~\ref{tab:bt-schema}, que se añaden sin alterar las existentes. De este modo, cualquier salida producida por BT-GRAPHRAG sigue siendo consumible por los módulos de búsqueda local y global del sistema base.

\begin{table}[h]
\centering
\small
\begin{tabular}{p{0.18\linewidth}p{0.22\linewidth}p{0.50\linewidth}}
\hline
\textbf{Tabla} & \textbf{Columna} & \textbf{Significado} \\ \hline
Relación & \texttt{t\_valid\_start} & Inicio del intervalo de validez en el mundo \\
Relación & \texttt{t\_valid\_end} & Fin del intervalo de validez (abierto si el hecho sigue siendo verdadero) \\
Relación & \texttt{t\_tx\_start} & Inicio del intervalo de creencia del sistema \\
Relación & \texttt{t\_tx\_end} & Fin del intervalo de creencia (abierto si el sistema sigue creyendo el hecho) \\
Relación & \texttt{cardinality} & Una de las cuatro categorías de cardinalidad \\
Relación & \texttt{status} & \texttt{active}, \texttt{disputed} o \texttt{retracted} \\
Relación & \texttt{support\_count} & Número de fuentes que corroboran la arista \\
Entidad  & \texttt{first\_seen} & Primera observación de la entidad en el corpus \\
Entidad  & \texttt{last\_seen} & Última observación \\
Entidad  & \texttt{active\_start} & Inicio inferido del período activo \\
Entidad  & \texttt{active\_end}   & Fin inferido del período activo (abierto si la entidad sigue activa) \\ \hline
\end{tabular}
\caption{Columnas que BT-GRAPHRAG añade a las tablas de relaciones y entidades del esquema de GraphRAG.}
\label{tab:bt-schema}
\end{table}

A diferencia de los sistemas analizados en la sección~\ref{sec:temporal-graphrag}, la temporalidad no actúa como un filtro aplicado a posteriori sobre un grafo estático, sino como una dimensión presente en cada etapa del \emph{pipeline} de indexación. Esta decisión de diseño resuelve la primera limitación identificada en el estado del arte y, al mismo tiempo, habilita el resto de los componentes: sin las marcas temporales en cada arista, ni los estados epistémicos de la sección~\ref{sec:bt-bitemporal} ni la búsqueda bidireccional de conflictos serían formulables.

%% ============================================================
\section{Modelo bitemporal a nivel de arista}\label{sec:bt-bitemporal}

\subsection{Cuádrupla de estado temporal}

Cada arista del grafo lleva una cuádrupla $(t_v^s, t_v^e, t_t^s, t_t^e)$ que codifica los dos ejes temporales descritos en la sección~\ref{sec:bitemporal}. Para una arista $a = (s, r, o)$, su estado temporal se define como
\begin{equation}
\mathrm{Q}(a) = \bigl(\,t_v^s,\; t_v^e,\; t_t^s,\; t_t^e\,\bigr) \;\in\; \mathcal{T}^4,
\end{equation}
donde $[t_v^s, t_v^e)$ es el \emph{tiempo válido} (intervalo en el que el hecho fue verdadero en el mundo) y $[t_t^s, t_t^e)$ es el \emph{tiempo de transacción} (intervalo en el que el sistema creyó el hecho). Para representar intervalos abiertos sin recurrir a valores nulos, BT-GRAPHRAG usa una constante $+\infty$:
\begin{equation}
t_v^e = +\infty \;\Longleftrightarrow\; \text{el hecho sigue siendo verdadero}, \qquad
t_t^e = +\infty \;\Longleftrightarrow\; \text{el sistema sigue creyéndolo}.
\end{equation}
Usar un valor explícito ($+\infty$) en lugar de un valor nulo simplifica las comparaciones temporales: todas las consultas se formulan como desigualdades sobre marcas temporales concretas, sin necesidad de ramas condicionales específicas para tratar los intervalos abiertos.

\subsection{Estados epistémicos derivados}

La cuádrupla temporal induce, mediante dos comparaciones contra $+\infty$, una partición del estado de cada arista en cuatro \emph{estados epistémicos} mutuamente excluyentes. Esta clasificación constituye la formalización operativa de los estados que la sección~\ref{sec:bitemporal} introdujo en el plano teórico.

\begin{itemize}
    \item \texttt{CURRENT\_TRUTH} ($t_v^e = +\infty \wedge t_t^e = +\infty$). Hecho actualmente verdadero y actualmente creído. Ejemplo: ``Tim Cook es CEO de Apple'' registrado en 2024 sin contradicción posterior.
    \item \texttt{HISTORICAL\_TRUTH} ($t_v^e < +\infty \wedge t_t^e = +\infty$). Hecho que fue verdadero en un intervalo cerrado del pasado. El sistema sigue creyendo que esa descripción del pasado es correcta. Ejemplo: ``Steve Jobs fue CEO de Apple entre 1997 y 2011''.
    \item \texttt{RETRACTED\_ERROR} ($t_v^e = +\infty \wedge t_t^e < +\infty$). Hecho que el sistema asumió inicialmente como verdadero pero que posteriormente fue retractado al identificarse como un error. Ejemplo: una extracción incorrecta del LLM que atribuyó un cargo a la persona equivocada y que un documento posterior desmiente sin aportar un hecho de reemplazo.
    \item \texttt{RETROACTIVE\_CORRECTION} ($t_v^e < +\infty \wedge t_t^e < +\infty$). Hecho histórico que el sistema mantuvo como verdadero durante un período de tiempo y que posteriormente fue desplazado por una versión corregida. Ejemplo: una arista ``Empresa X fundada en 1998'' que se sustituye por ``Empresa X fundada en 1999'' al disponer de una fuente de mayor fiabilidad. La arista original conserva su intervalo de transacción cerrado, mientras que la nueva queda registrada como la verdad vigente.
\end{itemize}

% TODO: figura fig:bt-cuadrupla-estados que represente la cuádrupla $(t_v^s, t_v^e, t_t^s, t_t^e)$ sobre dos ejes temporales (tiempo válido y tiempo de transacción) y muestre, en el mismo plano, los cuatro estados epistémicos derivados (CURRENT\_TRUTH, HISTORICAL\_TRUTH, RETRACTED\_ERROR, RETROACTIVE\_CORRECTION), cada uno con un ejemplo breve.
Estos estados no se almacenan como un campo independiente. Se derivan de la cuádrupla en tiempo de consulta. Esta derivación garantiza que el estado epistémico de cada arista nunca contradice sus marcas temporales: la única forma de moverla entre estados es modificar $t_v^e$ o $t_t^e$.

\subsection{Late arrivals y consultas as-of}

Un \emph{late arrival} es un documento cuya fecha del evento $t_{\text{event}}$ es bastante anterior a su fecha de incorporación al sistema $t_{\text{tx}}$. BT-GRAPHRAG marca como tales los documentos en los que la diferencia $t_{\text{tx}} - t_{\text{event}}$ supera un umbral. En estos casos, la búsqueda de conflictos no se ejecuta sobre el estado actual del grafo, sino sobre el estado que el grafo tenía en $t_{\text{event}}$. Una consulta \emph{as-of} a tiempo $T$ se define como
\begin{equation}
G_{\,\text{as-of}}(T) \;=\; \bigl\{\, a \in G \;\big|\; t_v^s(a) \le T < t_v^e(a) \;\wedge\; t_t^s(a) \le T_{\text{now}} < t_t^e(a) \,\bigr\},
\end{equation}
donde $T_{\text{now}}$ es el instante de la consulta. La primera condición restringe el grafo al tiempo válido consultado. La segunda mantiene solo las aristas que el sistema cree en el momento de la consulta. Esta combinación permite responder preguntas del tipo ``¿quién dirigía la empresa X en 2015, según lo que sabemos hoy?'' sin que las correcciones posteriores a 2015 contaminen la respuesta.

% TODO: figura fig:bt-as-of-query que muestre, sobre una línea temporal, una consulta as-of a tiempo $T$: el grafo completo arriba, las aristas activas en $T$ resaltadas abajo, y una arista de corrección retroactiva insertada con $t_t > T$ que la consulta debe excluir, ilustrando la doble condición sobre tiempo válido y tiempo de transacción.
Esta sección da respuesta a las limitaciones~4 y~6 identificadas en el estado del arte. La limitación~4 deriva de que ningún sistema de GraphRAG analizado integra la bitemporalidad en su \emph{pipeline}; en la propuesta presentada, la cuádrupla temporal forma parte de la propia tabla de aristas. La limitación~6 deriva de que los sistemas previos no distinguen formalmente entre la corrección de un hecho y su evolución natural; en BT-GRAPHRAG ambas situaciones se asocian a estados epistémicos bien diferenciados ---\texttt{RETROACTIVE\_CORRECTION} frente a \texttt{HISTORICAL\_TRUTH}--- que se infieren sin ambigüedad a partir de la cuádrupla.

%% ============================================================
\section{Clasificación de cardinalidad asimétrica}\label{sec:bt-cardinality}

\subsection{Las cuatro categorías}

La sección~\ref{sec:temporal-reasoning} estableció que la clasificación binaria exclusiva/no-exclusiva de TREK~\cite{trek2025} no captura la asimetría real de muchas relaciones. BT-GRAPHRAG sustituye esa clasificación binaria por una clasificación de cuatro categorías. Para cada tipo de relación $r$ con instancias $(s, r, o)$:

\begin{itemize}
    \item \texttt{SUBJECT\_EXCLUSIVE}: para un sujeto $s$ fijo, a lo sumo un objeto $o$ puede estar activo a la vez. Ejemplo: \texttt{HEADQUARTERED\_IN}. Una organización tiene una única sede activa, pero varias organizaciones pueden compartir la misma ciudad.
    \item \texttt{OBJECT\_EXCLUSIVE}: para un objeto $o$ fijo, a lo sumo un sujeto $s$ puede estar activo a la vez. Ejemplo: \texttt{IS\_CEO\_OF}. Una empresa tiene un único CEO activo, pero una persona puede ser CEO de varias empresas en paralelo.
    \item \texttt{BOTH\_EXCLUSIVE}: relación uno-a-uno; la exclusividad rige en ambos extremos. Ejemplo: \texttt{IS\_PRESIDENT\_OF}. Un país tiene un único presidente y, en la mayoría de los regímenes, una persona ocupa una única presidencia a la vez.
    \item \texttt{NON\_EXCLUSIVE}: ningún extremo impone exclusividad; las instancias pueden coexistir libremente. Ejemplos: \texttt{COLLABORATED\_WITH}, \texttt{PARTICIPATED\_IN}, \texttt{DEVELOPED}.
\end{itemize}

% TODO: figura fig:bt-cardinalidad-cuatro-categorias que ilustre, mediante diagramas de aristas con sujeto y objeto fijados o libres, las cuatro categorías de cardinalidad (SUBJECT\_EXCLUSIVE, OBJECT\_EXCLUSIVE, BOTH\_EXCLUSIVE, NON\_EXCLUSIVE) acompañadas de un ejemplo prototípico para cada una (HEADQUARTERED\_IN, IS\_CEO\_OF, IS\_PRESIDENT\_OF, COLLABORATED\_WITH).
La utilidad práctica de esta distinción se aprecia con claridad en el ejemplo \texttt{nació\_en}, discutido en la sección~\ref{sec:temporal-reasoning}. Desde la perspectiva del sujeto, la relación es exclusiva: una persona nace en un único lugar. Desde la perspectiva del objeto, deja de serlo: en una misma ciudad pueden nacer muchas personas. En el esquema de TREK, clasificar \texttt{nació\_en} como ``exclusiva'' generaría falsos conflictos cada vez que dos personas distintas nacieran en la misma ciudad; clasificarla como ``no exclusiva'', en cambio, dejaría pasar el conflicto real cuando dos documentos atribuyeran a una misma persona dos lugares de nacimiento diferentes. En BT-GRAPHRAG, la relación se clasifica como \texttt{SUBJECT\_EXCLUSIVE}, de modo que la búsqueda de conflictos se activa únicamente desde el sujeto, capturando la asimetría sin introducir ruido.

\subsection{Clasificación por LLM en tiempo de indexación}

La clasificación se realiza una única vez por tipo de relación, no por instancia. Tras la consolidación de tipos descrita en la sección~\ref{sec:bt-cger}, cada tipo nuevo se envía a un LLM mediante un \emph{prompt} que incluye ejemplos del grafo y las definiciones de las cuatro categorías. El modelo devuelve exactamente una etiqueta. El \emph{prompt} favorece por defecto la categoría \texttt{NON\_EXCLUSIVE}: solo asigna una categoría exclusiva cuando existe una restricción del mundo real claramente establecida. Esta política reduce los falsos positivos.

La etiqueta resultante se almacena como atributo del tipo de relación. El sistema admite además un mapa de sobreescrituras manual para los casos en que un experto de dominio desee forzar una clasificación concreta. Dado que la clasificación opera a nivel de tipo y no de arista, su costo es reducido incluso en corpus de gran tamaño.

\subsection{Dirección de la búsqueda de conflictos}

La cardinalidad asociada a una arista determina la naturaleza de las consultas que el sistema realiza sobre el grafo para detectar conflictos, proceso que se describe en detalle en la sección~\ref{sec:bt-etcdr}. Para una arista candidata $(s, r, o)$, el comportamiento es el siguiente:

\begin{itemize}
    \item Si la cardinalidad de $r$ es \texttt{SUBJECT\_EXCLUSIVE}, el sistema busca aristas activas de la forma $(s, r, *)$ con objeto distinto, es decir, otros objetos vinculados al mismo sujeto mediante el mismo tipo de relación.
    \item Si es \texttt{OBJECT\_EXCLUSIVE}, busca aristas activas de la forma $(*, r, o)$ con sujeto distinto. Esta es precisamente la consulta que ZEP/Graphiti~\cite{zep2025} omite, tal como se discute en la sección~\ref{sec:bt-etcdr}.
    \item Si es \texttt{BOTH\_EXCLUSIVE}, se ejecutan ambas consultas de forma simultánea.
    \item Si es \texttt{NON\_EXCLUSIVE}, se omiten ambas consultas estructurales y el sistema se limita a comprobar si la pareja $(s, o)$ ya dispone de una arista del mismo tipo, con el fin de detectar duplicados por corroboración.
\end{itemize}

% TODO: figura fig:bt-cardinalidad-busqueda que muestre, para una arista candidata $(s, r, o)$, qué consulta(s) sobre el grafo se activan según la cardinalidad de $r$: flechas desde $s$ (subject-side) y desde $o$ (object-side), con las cuatro categorías indicando cuál(es) de las dos consultas se ejecuta(n).
Esta sección da respuesta a la limitación~3 identificada en el estado del arte. La cardinalidad deja de ser una etiqueta simétrica de uso interno, como ocurre en TREK, para convertirse en un mecanismo de despacho que gobierna la dirección de la búsqueda de conflictos en el siguiente módulo del \emph{pipeline}.

%% ============================================================
\section{Resolución contextual y temporal de entidades (CGER y CGRR)}\label{sec:bt-cger}

CGER y CGRR son dos módulos incorporados al \emph{pipeline} de GraphRAG con el propósito de abordar la resolución de entidades y de tipos de relación desde una perspectiva contextual y temporal. CGER (\emph{Cross-Graph Entity Resolution}) sustituye la coincidencia exacta de nombres del GraphRAG original por un procedimiento en dos fases: un \emph{scoring} basado en la similitud coseno entre los \emph{embeddings} de descripción de las entidades y, para los pares cuyo coseno supera un umbral, una verificación por LLM que razona conjuntamente sobre las descripciones y la temporalidad de cada entidad. CGRR (\emph{Cross-Graph Relationship Resolution}) aplica una idea análoga a los tipos de relación, con un \emph{scoring} compuesto adaptado a la naturaleza de los predicados.

\subsection{Pre-filtrado y \emph{scoring} por similitud de descripción}

La comparación exhaustiva de cada entidad nueva contra la totalidad de entidades presentes en el grafo resulta computacionalmente costosa en corpus de gran tamaño. Para acotar el espacio de búsqueda, CGER realiza una consulta por similitud sobre los \emph{embeddings} de descripción de las entidades: dado un candidato, se recuperan sus $K$ vecinos más próximos según la similitud coseno~\cite{TODO_cosine_similarity}. El parámetro $K$ es configurable en función del dominio y del tamaño del corpus. Sobre este subconjunto, restringido además a los candidatos del mismo tipo, CGER calcula la similitud coseno entre los \emph{embeddings} de descripción de la entidad nueva y los del candidato, y conserva como mejor coincidencia aquella con mayor coseno. Esta operación, lineal respecto al número de entidades nuevas, constituye la única señal de \emph{scoring} del módulo: el resto de las decisiones se delegan al LLM, que dispone de información cualitativa más rica que cualquier función agregada de señales heterogéneas.

\subsection{Verificación por LLM con contexto temporal}\label{subsec:bt-cger-llm}

La similitud coseno entre descripciones es un indicador robusto de proximidad semántica pero, por sí sola, no distingue entre dos referentes distintos descritos en términos parecidos y dos estados sucesivos del mismo referente. El ejemplo ``Apple 1985 vs.\ 2024'', introducido en la sección~\ref{subsec:entity-resolution}, ilustra esta limitación: las descripciones, los \emph{embeddings} y la propia similitud coseno coinciden de forma casi perfecta entre ambas menciones, y sin embargo fusionarlas en un único nodo eliminaría la diferencia de estado entre los dos períodos. Para resolver estos casos, todo par con coseno por encima del umbral $\tau_{\cos}$ se somete a una verificación por LLM. El \emph{prompt} entrega al modelo, para cada una de las dos entidades, su nombre, su tipo, su descripción y el intervalo $[\,\texttt{active\_start}, \texttt{active\_end}\,]$ que representa su período activo, y le solicita exactamente una de tres etiquetas:

\begin{itemize}
    \item \texttt{SAME}: las dos entidades se refieren al mismo referente y sus períodos activos son lo suficientemente próximos como para que la fusión no elimine información de estado. Es el único veredicto que produce un merge.
    \item \texttt{DIFFERENT\_ENTITY}: las entidades se refieren a referentes distintos del mundo real (por ejemplo, una entidad genérica frente a una específica, dos ediciones de un mismo evento o dos organizaciones con tokens compartidos). Se preservan como nodos diferenciados.
    \item \texttt{DIFFERENT\_TEMPORAL}: las entidades describen al mismo referente pero la separación temporal entre sus períodos activos es lo bastante grande como para que su fusión destruya estados claramente distintos del referente. Este veredicto materializa el caso ``Apple 1985 vs.\ 2024'' y se aplica también a configuraciones más sutiles —como dos plantillas sucesivas de un mismo equipo deportivo o dos gabinetes consecutivos de una misma cartera ministerial—. Las entidades se preservan como nodos diferenciados.
\end{itemize}

El umbral $\tau_{\cos}$ controla el equilibrio entre coste y cobertura: cuanto más bajo se sitúe, más pares se someten a verificación por LLM (incrementando la cobertura, a costa de un mayor número de llamadas), y cuanto más alto, menos llamadas se efectúan, asumiendo el riesgo de no detectar fusiones legítimas cuyas descripciones difieran ligeramente. La decisión sobre la temporalidad recae íntegramente sobre el LLM, que es el único componente del módulo con capacidad de razonar conjuntamente sobre las descripciones y los períodos activos. Esta delegación es deliberada: la temporalidad es una señal cuya \emph{relevancia} depende del referente —décadas son significativas para una organización, meses lo son para una plantilla deportiva, días lo son para un gabinete de gobierno—, y cualquier combinación lineal entre solapamiento temporal y similitud semántica obliga a fijar a priori una sensibilidad común a todos los tipos de entidad. El LLM, en cambio, modula su criterio en función del tipo y del contenido de la descripción.

% TODO: figura fig:bt-cger-pipeline que represente las dos fases de CGER como un pipeline horizontal: (1) prefiltrado por similitud vectorial (top-K) sobre el embedding de descripción; (2) verificación por LLM con contexto temporal, que recibe nombre, tipo, descripción y período activo de cada entidad y devuelve uno de los tres veredictos SAME / DIFFERENT_ENTITY / DIFFERENT_TEMPORAL.

\subsection{Resolución de tipos de relación}

CGRR aplica una idea análoga a los tipos de relación, pero su \emph{scoring} se mantiene compuesto: a diferencia de las entidades, los tipos de relación no portan un período activo propio, por lo que la dimensión temporal queda fuera de juego. La función de \emph{scoring} combina BM25 sobre el nombre del tipo, similitud semántica de la descripción y un indicador de coincidencia de extremos: cuando los pares sujeto-objeto se solapan, la evidencia de equivalencia semántica se refuerza. Así, por ejemplo, el sistema puede fusionar \texttt{IS\_CEO\_OF} y \texttt{LEADS} cuando el contexto indica que ambas etiquetas designan el mismo predicado. Este proceso de consolidación reduce la dispersión del vocabulario de relaciones antes de invocar al módulo de clasificación de cardinalidad y al de detección de conflictos.

En esta sección se da respuesta a la limitación~5 identificada en el estado del arte. La resolución de entidades en sistemas como DALK y DYNAMO es de naturaleza puramente semántica y, como se señala en la sección~\ref{subsec:entity-resolution}, no contempla la posibilidad de que dos menciones pertenecientes a períodos distintos describan estados sustancialmente diferentes de una misma entidad. CGER aborda esta carencia con una arquitectura que separa explícitamente las dos preguntas implicadas: ¿son las descripciones lo bastante próximas para sospechar identidad? (responde la similitud coseno, como pre-filtro); y, si lo son, ¿debemos fusionar las dos menciones? (responde el LLM, considerando descripción y temporalidad de manera conjunta, y distinguiendo entre referentes distintos y estados temporalmente separados de un mismo referente). Por su parte, la consolidación de tipos que aporta CGRR garantiza que tanto la clasificación de cardinalidad como la detección de conflictos operen sobre un vocabulario canónico y sin ambigüedades léxicas.

%% ============================================================
\section{Detección y resolución de conflictos temporales (ETCDR)}\label{sec:bt-etcdr}

ETCDR (\emph{Edge-level Temporal Conflict Detection and Resolution}) constituye el módulo central del sistema y el que le otorga nombre. A diferencia de T-GRAG~\cite{tgrag2025} y TG-RAG~\cite{tg_rag2025}, no delega la coherencia temporal a la fase de consulta, sino que detecta y resuelve los conflictos durante la propia indexación, dejando constancia explícita de cada decisión adoptada.

\subsection{Detección bidireccional dirigida por cardinalidad}

Para cada arista candidata $(s, r, o)$, el sistema ejecuta como máximo tres familias de consultas, seleccionadas en función de la cardinalidad de $r$:

\begin{itemize}
    \item \textbf{Consulta desde el sujeto} (relaciones \texttt{SUBJECT\_EXCLUSIVE} o \texttt{BOTH\_EXCLUSIVE}). Recupera las aristas activas de la forma $(s, r, o')$ con $o' \neq o$. Detecta conflictos del tipo ``el sujeto ya tiene asignado un objeto distinto para este tipo de relación''.
    \item \textbf{Consulta desde el objeto} (relaciones \texttt{OBJECT\_EXCLUSIVE} o \texttt{BOTH\_EXCLUSIVE}). Recupera las aristas activas de la forma $(s', r, o)$ con $s' \neq s$. Detecta conflictos del tipo ``el objeto ya tiene asignado un sujeto distinto para este tipo de relación''.
    \item \textbf{Consulta sobre la misma pareja} (todas las cardinalidades). Recupera las aristas activas de la forma $(s, r, o)$ con el mismo sujeto y objeto, con el fin de identificar candidatos que deben ser corroborados.
\end{itemize}

La consulta desde el objeto es precisamente la que ZEP/Graphiti~\cite{zep2025} omite. Como se señala en la sección~\ref{sec:bitemporal}, este sistema es el que más se aproxima al modelo bitemporal entre los analizados en el estado del arte; sin embargo, detecta contradicciones únicamente por similitud semántica entre aristas que comparten sujeto, lo que cubre el conflicto desde el lado del sujeto, pero no desde el del objeto. En consecuencia, si dos personas distintas se atribuyen el mismo cargo en la misma empresa durante el mismo intervalo de tiempo, el sistema no genera ninguna alerta. ETCDR subsana esta omisión: las categorías \texttt{OBJECT\_EXCLUSIVE} y \texttt{BOTH\_EXCLUSIVE} activan explícitamente la consulta desde el objeto, garantizando la detección de este tipo de conflicto.

Las consultas se ejecutan tanto sobre el grafo ya persistido como sobre el lote de aristas en memoria pendientes de escritura. Este doble alcance es necesario porque dos aristas extraídas en el mismo lote pueden entrar en conflicto entre sí antes de haber sido incorporadas al grafo.

\subsection{Decision Router}

Cuando alguna de las consultas devuelve aristas conflictivas, ETCDR delega la decisión a un componente basado en LLM denominado \emph{Decision Router}. Este recibe como entrada: (i) la arista existente con su cuádrupla y descripción, (ii) la arista candidata con su cuádrupla y descripción, (iii) la cardinalidad del tipo de relación, (iv) el tipo de conflicto detectado y (v) la procedencia y el grado de confianza asociado a cada extremo. La salida del componente es una de las cinco estrategias descritas en el subapartado siguiente, acompañada de un valor de confianza en el intervalo $[0, 1]$ y una breve justificación. Cuando dicho valor queda por debajo de un umbral configurable, el sistema degrada automáticamente la decisión a \texttt{DISAGREEMENT}, evitando así que resoluciones de fiabilidad insuficiente se propaguen al grafo.

\begin{algorithm}[H]
\caption{Resolución de un candidato en ETCDR}\label{alg:etcdr}
\KwIn{Arista candidata $a = (s, r, o)$ con cuádrupla $\widetilde{Q}$}
$c \gets \texttt{cardinality}(r)$\;
$C_S \gets \emptyset$;\, $C_O \gets \emptyset$;\, $C_P \gets \emptyset$\;
\If{$c \in \{\texttt{SUBJECT\_EXCLUSIVE}, \texttt{BOTH\_EXCLUSIVE}\}$}{
    $C_S \gets$ aristas activas $(s, r, *)$ con objeto $\neq o$\;
}
\If{$c \in \{\texttt{OBJECT\_EXCLUSIVE}, \texttt{BOTH\_EXCLUSIVE}\}$}{
    $C_O \gets$ aristas activas $(*, r, o)$ con sujeto $\neq s$\;
}
$C_P \gets$ aristas activas $(s, r, o)$\;
\If{$C_S \cup C_O \cup C_P = \emptyset$}{
    \Return \texttt{NEW\_EDGE}\;
}
$(\sigma, \kappa) \gets \texttt{DecisionRouterLLM}(a, C_S \cup C_O \cup C_P, c)$\;
\If{$\kappa < \tau_{\text{conf}}$}{
    \Return \texttt{DISAGREEMENT}\;
}
\Return $\sigma$\;
\end{algorithm}

% TODO: figura fig:bt-etcdr-flowchart que diagrame el flujo del Algoritmo~\ref{alg:etcdr}: arista candidata $\to$ consulta de cardinalidad $\to$ tres familias de consultas (subject-side, object-side, same-pair) condicionadas por la categoría $\to$ Decision Router (LLM) $\to$ ramificación a las cinco estrategias (EVOLUTION, CORRECTION, CORROBORATION, DISAGREEMENT, NEW\_EDGE), con el umbral de confianza que degrada a DISAGREEMENT.
\subsection{Cinco estrategias de resolución}

Las cinco estrategias de resolución y su efecto sobre la cuádrupla se recogen en la Tabla~\ref{tab:etcdr-strategies}. Desde el punto de vista operativo, cualquier resolución se reduce a, como máximo, dos modificaciones sobre los campos \texttt{t\_valid\_end} o \texttt{t\_tx\_end} y la inserción opcional de una nueva arista. El hecho de que toda decisión de resolución se exprese como un cambio acotado sobre la cuádrupla es lo que garantiza tanto la trazabilidad completa como la preservación del estado anterior. Las aristas históricas se cierran, no se eliminan.

\begin{table}[h]
\centering
\small
\begin{tabular}{p{0.16\linewidth}p{0.24\linewidth}p{0.18\linewidth}p{0.16\linewidth}p{0.20\linewidth}}
\hline
\textbf{Estrategia} & \textbf{Condición} & \textbf{Efecto sobre $t_v^e$ existente} & \textbf{Efecto sobre $t_t^e$ existente} & \textbf{Otros efectos} \\ \hline
\texttt{EVOLUTION} & El mundo cambió genuinamente & Se cierra al $t_v^s$ del candidato & Sin cambio & Inserta candidato como \texttt{CURRENT\_TRUTH} \\
\texttt{CORRECTION} & La arista existente era errónea & Sin cambio & Se cierra a $t_{\text{now}}$ & Inserta candidato heredando el $t_v$ del original \\
\texttt{CORROBORATION} & Candidato confirma la existente & Sin cambio & Sin cambio & Incrementa \texttt{support\_count} y agrega procedencia \\
\texttt{DISAGREEMENT} & Ambigüedad o evidencia conflictiva & Sin cambio & Sin cambio & Inserta candidato con \texttt{status = disputed} \\
\texttt{NEW\_EDGE} & Sin conflicto detectado & --- & --- & Inserta candidato como arista independiente \\ \hline
\end{tabular}
\caption{Estrategias de resolución de ETCDR y su efecto sobre la cuádrupla bitemporal de la arista existente.}
\label{tab:etcdr-strategies}
\end{table}

% TODO: figura fig:bt-etcdr-strategies-effect que muestre, sobre el plano (tiempo válido, tiempo de transacción), el efecto de cada una de las cinco estrategias sobre la cuádrupla de la arista existente y, en su caso, la inserción de una nueva arista; debe quedar visualmente clara la diferencia entre EVOLUTION (cierra $t_v^e$) y CORRECTION (cierra $t_t^e$).
La distinción entre \texttt{EVOLUTION} y \texttt{CORRECTION} constituye la traducción operativa de la diferencia conceptual entre evolución natural y corrección retroactiva, introducidas en las secciones~\ref{sec:graphrag-foundations} y~\ref{sec:bitemporal}, respectivamente. \texttt{EVOLUTION} cierra el tiempo válido de la arista anterior porque el hecho fue verdadero hasta ese momento, instante en el que fue desplazado por uno nuevo. \texttt{CORRECTION}, en cambio, cierra el tiempo de transacción, dado que el sistema deja de creer en la arista anterior al reconocer que el hecho subyacente nunca fue verdadero. \texttt{CORROBORATION} no modifica la cuádrupla, pero refuerza la confianza en la arista existente mediante el incremento de \texttt{support\_count}. Por último, \texttt{DISAGREEMENT} preserva ambas versiones marcando la nueva como \texttt{disputed}, lo que permite al módulo de consulta, expuesto en la sección~\ref{sec:bt-integration}, tratarla de forma diferenciada.

\subsection{Trazabilidad de las resoluciones}

Cada decisión adoptada por ETCDR se serializa en un registro estructurado de auditoría. Cada entrada contiene la arista candidata, las aristas conflictivas detectadas, el tipo de conflicto, la cardinalidad consultada, la estrategia seleccionada, el valor de confianza emitido por el \emph{Decision Router} y la justificación textual del modelo. Este registro cumple una doble función. Como herramienta de auditoría, permite reconstruir el razonamiento que llevó a marcar una arista como \texttt{disputed} o a retraer otra. Como instrumento de evaluación, sustenta los experimentos del capítulo siguiente: contrastar la calidad de las decisiones de ETCDR con anotaciones manuales requiere precisamente este nivel de granularidad.

A diferencia de T-GRAG y TG-RAG (sección~\ref{sec:temporal-graphrag}), que delegan toda la coherencia temporal al filtrado por fecha en la fase de consulta y no detectan ni resuelven conflictos durante la indexación, ETCDR los aborda en el mismo \emph{pipeline} en el que se incorpora el conocimiento al grafo. La consecuencia directa es que el grafo resultante no contiene contradicciones latentes a la espera de ser enmascaradas por una restricción temporal impuesta por el usuario. Las únicas contradicciones presentes son aquellas que el sistema ha clasificado explícitamente como \texttt{DISAGREEMENT}, de las cuales queda constancia auditable en el registro de decisiones.

%% ============================================================
\section{Integración con el pipeline de GraphRAG y consulta temporal}\label{sec:bt-integration}

\subsection{Reportes de comunidad temporales}

La jerarquía (pirámide) de comunidades de Leiden generada por GraphRAG~\cite{edge2024local} se conserva íntegramente. BT-GRAPHRAG no modifica el algoritmo de partición ni la estructura jerárquica que produce; la intervención se concentra exclusivamente en la generación de los reportes de comunidad. Para cada nivel de la jerarquía y cada comunidad, el \emph{prompt} de generación del reporte se reemplaza por una variante enriquecida que recibe, además de las entidades y relaciones del subgrafo, las marcas temporales de cada arista. La salida incorpora una sección de \emph{evolución temporal} a las secciones habituales ---título, resumen, hallazgos y valoración---. Dicha sección sintetiza, en lenguaje natural, las transiciones relevantes del período analizado: la aparición o desaparición de entidades, los cambios de estado en las relaciones y los hechos que se encuentran actualmente en disputa.

% TODO: figura fig:bt-temporal-report que muestre un fragmento de reporte temporal con sus secciones (resumen, evolución, hallazgos) marcadas.

Los reportes enriquecidos mantienen el formato esperado por el módulo de búsqueda global de GraphRAG, por lo que dicho módulo opera sobre ellos sin requerir modificación alguna. Las respuestas generadas incorporan narrativa temporal de forma natural, como reflejo directo del contenido de los reportes. La búsqueda local se beneficia igualmente de esta integración: al recuperar el vecindario de las entidades semilla, las descripciones de las aristas ya portan las marcas temporales, lo que permite al modelo situarlas en su período cronológico correcto.

\subsection{Descomposición temporal de consultas}

Las consultas que combinan referencias a múltiples instantes o períodos de tiempo ---por ejemplo, ``¿quién dirigía la empresa X cuando se fusionó con Y?''--- no pueden resolverse mediante un único filtro temporal. BT-GRAPHRAG las descompone en sub-consultas independientes, cada una asociada a una restricción temporal explícita. Para cada sub-consulta, el sistema determina su tipo ---\texttt{POINT\_IN\_TIME}, \texttt{RANGE}, \texttt{EVOLUTION} o \texttt{COMPARISON}--- y la ejecuta sobre el grafo en el modo \emph{as-of} correspondiente. Los resultados parciales se agregan mediante una síntesis temporal que hace explícitas las transiciones entre períodos. Las aristas marcadas como \texttt{disputed} se someten a un proceso de resolución de disputas que pondera la fecha de cada versión, la confianza de la fuente y el valor de \texttt{support\_count}.

La descomposición temporal de consultas no es un enfoque novedoso en sí mismo: T-GRAG~\cite{tgrag2025} y MRAG~\cite{mrag2025} introdujeron variantes de esta misma idea. La distinción fundamental de BT-GRAPHRAG reside en que dicha descomposición opera sobre un grafo en el que la coherencia temporal ha sido garantizada durante la indexación: los conflictos están resueltos o explícitamente etiquetados, las entidades se encuentran temporalmente diferenciadas y la cardinalidad es asimétrica. En consecuencia, el filtrado \emph{as-of} aplicado a cada sub-consulta devuelve un subgrafo internamente coherente, sin necesidad de reconciliar contradicciones en tiempo de consulta.

\subsection{Actualización incremental del corpus}

La incorporación incremental de nuevos documentos constituye un escenario natural de operación para BT-GRAPHRAG. Cada arista nueva atraviesa las mismas siete etapas del \emph{pipeline} y se reconcilia con el estado existente del grafo mediante los módulos CGER, CGRR y ETCDR. Para evitar el reprocesamiento de documentos ya indexados, el sistema mantiene un registro persistente de los documentos previamente incorporados al grafo, de modo que únicamente los documentos nuevos son admitidos en cada lote de procesamiento.

% TODO: figura fig:bt-incremental-update que ilustre el flujo de actualización incremental: documentos nuevos $\to$ entidades modificadas $\to$ expansión $k$-hop del vecindario $\to$ comunidades obsoletas marcadas $\to$ regeneración selectiva de reportes; las comunidades no afectadas deben aparecer claramente preservadas sin modificación.
La regeneración de los reportes de comunidad sigue el patrón propuesto por TG-RAG~\cite{tg_rag2025}. En lugar de reconstruir la pirámide en su totalidad, BT-GRAPHRAG identifica los nodos afectados por las aristas nuevas o modificadas, calcula su vecindario hasta una distancia $k$ configurable y reejecuta la partición de Leiden y la generación de reportes exclusivamente sobre las comunidades que intersecan dicho vecindario. Las comunidades no afectadas conservan sus reportes previos sin modificación. Esta política reduce el coste amortizado de la actualización a una fracción del coste de la indexación inicial, lo que hace viable mantener el grafo actualizado en escenarios de ingestión continua.

%% ============================================================
\section{Conclusiones}\label{sec:bt-conclusions}

El presente capítulo ha descrito el diseño conceptual de BT-GRAPHRAG, una extensión del \emph{pipeline} de GraphRAG que integra semántica bitemporal a nivel de arista y resuelve los conflictos temporales durante la fase de indexación, en lugar de trasladar su gestión a la fase de consulta. La propuesta se articula en torno a cuatro componentes principales:

\begin{itemize}
    \item la cuádrupla temporal y los cuatro estados epistémicos que de ella se derivan,
    \item la clasificación asimétrica de cardinalidad en cuatro categorías,
    \item los módulos CGER y CGRR de resolución contextual y temporal de entidades y tipos de relación, y
    \item el módulo ETCDR de detección bidireccional de conflictos y resolución mediante \emph{Decision Router}.
\end{itemize}

Estos componentes se insertan en un \emph{pipeline} de siete etapas que preserva la compatibilidad con los modos de búsqueda local, global y DRIFT del sistema base, al tiempo que mantiene una representación bitemporal completa capaz de sostener las consultas \emph{as-of} y la detección bidireccional dirigida por cardinalidad.

La Tabla~\ref{tab:bt-vs-soa} extiende la Tabla~\ref{tab:soa-comparison} del estado del arte y recoge, para cada una de las seis limitaciones identificadas, el componente de BT-GRAPHRAG que la resuelve.

\begin{table}[h]
\centering
\small
\begin{tabular}{p{0.06\linewidth}p{0.36\linewidth}p{0.50\linewidth}}
\hline
\textbf{Lim.} & \textbf{Limitación del SoA (\ref{sec:soa-conclusions})} & \textbf{Respuesta de BT-GRAPHRAG} \\ \hline
1 & Temporalidad presente pero no integrada en el \emph{pipeline} de GraphRAG. & \emph{Pipeline} de siete etapas con la temporalidad como atributo de primera clase y representación dual del grafo (\ref{sec:bt-pipeline}). \\
2 & Separación temporal en lugar de resolución de conflictos. & ETCDR: detección bidireccional dirigida por cardinalidad, \emph{Decision Router} y cinco estrategias de resolución (\ref{sec:bt-etcdr}). \\
3 & Cardinalidad simétrica de TREK (exclusiva/no-exclusiva). & Clasificación de cuatro categorías \texttt{SUBJECT\_EXCLUSIVE}, \texttt{OBJECT\_EXCLUSIVE}, \texttt{BOTH\_EXCLUSIVE}, \texttt{NON\_EXCLUSIVE} (\ref{sec:bt-cardinality}). \\
4 & Bitemporalidad fuera del \emph{pipeline} de GraphRAG. & Cuádrupla $(t_v^s, t_v^e, t_t^s, t_t^e)$ embebida en cada arista del modelo de datos (\ref{sec:bt-bitemporal}). \\
5 & Resolución de entidades sin temporalidad. & CGER con pre-filtro por similitud coseno de descripciones y verificación LLM con contexto temporal y veredicto \texttt{DIFFERENT\_TEMPORAL} (\ref{sec:bt-cger}). \\
6 & Corrección y evolución sin distinción formal. & Estados epistémicos \texttt{HISTORICAL\_TRUTH} vs.\ \texttt{RETROACTIVE\_CORRECTION} y estrategias \texttt{EVOLUTION} vs.\ \texttt{CORRECTION} en ETCDR (\ref{sec:bt-bitemporal}, \ref{sec:bt-etcdr}). \\ \hline
\end{tabular}
\caption{Correspondencia uno-a-uno entre las seis limitaciones del estado del arte y los componentes de BT-GRAPHRAG que las resuelven.}
\label{tab:bt-vs-soa}
\end{table}

La correspondencia recogida en la Tabla~\ref{tab:bt-vs-soa} evidencia que el diseño no constituye una mera agregación de mecanismos independientes, sino una arquitectura cohesionada en la que cada componente habilita al siguiente: la cuádrupla temporal fundamenta los estados epistémicos; los estados epistémicos se materializan a través de las cinco estrategias de resolución de ETCDR; este se dirige por la clasificación asimétrica de cardinalidad; dicha clasificación opera sobre tipos previamente consolidados por CGRR; y CGER delega en el LLM la decisión de fusión sobre los pares con coseno alto, suministrándole los períodos activos de cada entidad para que pueda distinguir entre referentes distintos y estados temporalmente separados del mismo referente. La interdependencia entre componentes es, por tanto, estructural: ninguno de ellos tiene sentido de forma aislada.

El capítulo siguiente aborda la evaluación experimental de la propuesta. La estrategia de evaluación cubre tres dimensiones: la corrección temporal de las respuestas a consultas con restricciones \emph{as-of}, la calidad de las decisiones adoptadas por ETCDR frente a anotaciones manuales sobre conflictos sembrados en el corpus, y el coste amortizado de la actualización incremental sobre corpus que evolucionan a lo largo del tiempo.
